\section{Relative Spec}
\label{mk-chap:Picard_RelativeSpec}

\section{Quasi-coherent sheaves of algebras}
\label{mk-sec:qc_sheaf_of_algebras}

\begin{definition}[Quasi-coherent \(\mathcal{O}_X\)-algebra]
  \label{mk-def:qc_sheaf_of_algebras}
  \dcref{01LL,custom}
  \textit{Source: \cite{stacks-project}, tag 01LL (situation-relative-spec).}
  Let \(X\) be a scheme. A \emph{quasi-coherent sheaf of \(\mathcal{O}_X\)-algebras} is a
  sheaf \(\mathcal{A}\) of \(\mathcal{O}_X\)-algebras (i.e. an \(\mathcal{O}_X\)-module
  equipped with \(\mathcal{O}_X\)-linear multiplication and unit morphisms satisfying
  associativity, commutativity, and the unit law as sheaf morphisms) whose underlying
  \(\mathcal{O}_X\)-module is quasi-coherent. We write \(\mathrm{QcohAlg}(X)\) for the
  category of quasi-coherent sheaves of \(\mathcal{O}_X\)-algebras and
  \(\mathcal{O}_X\)-algebra morphisms.
\end{definition}

\section{The relative-spectrum scheme}
\label{mk-sec:relspec_scheme}

The construction is given affine-locally: on every affine open \(U = \Spec R \subseteq X\) one
sets \(\pi^{-1}(U) := \Spec(\mathcal{A}(U))\) regarded as an \(R\)-algebra, and the local pieces
glue compatibly because \(\mathcal{A}\) is quasi-coherent (the transition isomorphism
\(\mathcal{A}(U) \otimes_R S \xrightarrow{\sim} \mathcal{A}(V)\) for an inclusion
\(V = \Spec S \subseteq U\) gives an open immersion of the corresponding Specs).

\begin{theorem}[Relative spectrum exists]
  \label{mk-thm:relative_spec_exists}
  \dcref{01LQ,custom}
  \uses{mk-def:qc_sheaf_of_algebras}
  \textit{Source: \cite{stacks-project}, tag 01LQ (lemma-glue-relative-spec);
  cf.\ \cite{hartshorne-ag}, II~Ex.~5.17(a).}
  Let \(X\) be a scheme and \(\mathcal{A}\) a quasi-coherent \(\mathcal{O}_X\)-algebra
  (\ref{mk-def:qc_sheaf_of_algebras}). There exists a scheme
  \(\underline{\Spec}_X(\mathcal{A})\) together with an affine morphism
  \(\pi : \underline{\Spec}_X(\mathcal{A}) \to X\) such that
  \begin{enumerate}
    \item for every affine open \(U \subseteq X\), there is an isomorphism
      \(i_U : \pi^{-1}(U) \xrightarrow{\sim} \Spec(\mathcal{A}(U))\) over \(U\);
    \item for \(U \subseteq U' \subseteq X\) both affine open, the composite
      $\Spec(\mathcal{A}(U)) \xrightarrow{i_U^{-1}} \pi^{-1}(U) \hookrightarrow
      \pi^{-1}(U') \xrightarrow{i_{U'}} \Spec(\mathcal{A}(U'))$
      coincides with the open immersion induced by the restriction
      \(\mathcal{A}(U') \to \mathcal{A}(U)\).
  \end{enumerate}
  The pair \((\underline{\Spec}_X(\mathcal{A}), \pi)\) is unique up to unique isomorphism
  over \(X\).
\end{theorem}

\begin{proof}
  By the open-cover gluing principle for schemes, it suffices to specify affine local
  pieces \(\Spec(\mathcal{A}(U))\) for each affine open \(U \subseteq X\) and
  open-immersion transition data
  \(\Spec(\mathcal{A}(U)) \hookrightarrow \Spec(\mathcal{A}(U'))\) for \(U \subseteq U'\)
  satisfying the cocycle condition. The transition morphisms come from the ring maps
  \(\mathcal{A}(U') \to \mathcal{A}(U)\); quasi-coherence yields the cartesian square
  \(\Spec(\mathcal{A}(U)) = U \times_{U'} \Spec(\mathcal{A}(U'))\), so each transition is an
  open immersion. Transitivity (Stacks lemma-transitive-spec) gives the cocycle condition.
  Uniqueness of the glued scheme is the standard property of the gluing functor.
\end{proof}

\begin{definition}[Structure morphism of the relative spectrum]
  \label{mk-def:relspec_structure_morphism}
  \dcref{custom}

  \uses{mk-thm:relative_spec_exists, mk-def:qc_sheaf_of_algebras}
  Let \(X\) be a scheme and \(\mathcal{A}\) a quasi-coherent \(\mathcal{O}_X\)-algebra.
  The \emph{structure morphism}
  \(\pi : \underline{\Spec}_X(\mathcal{A}) \to X\) is the canonical projection of the
  relative spectrum to its base, obtained as the colimit descent of the affine-local
  structure morphisms \(\Spec(\mathcal{A}(U)) \to U\) over the directed affine open
  cover of \(X\). It is the morphism whose affineness and base-change behaviour are the
  subject of the universal-property and base-change results below.
\end{definition}

\begin{theorem}[Universal property of the relative spectrum]
  \label{mk-thm:relative_spec_univ}
  \dcref{01LQ,custom}
  \uses{mk-thm:relative_spec_exists, mk-thm:relative_spec_affine_base,
    mk-def:relspec_structure_morphism}
  \textit{Source: \cite{stacks-project}, tag 01LQ (lemma-spec + section-spec
  functor description); \cite{hartshorne-ag}, II~Ex.~5.17(b).}
  Let \(X\) be a scheme and \(\mathcal{A}\) a quasi-coherent \(\mathcal{O}_X\)-algebra.
  Then \(\underline{\Spec}_X(\mathcal{A})\) represents the functor
  \(F : \Sch^{op} \to \mathrm{Set}\),
  $T \mapsto F(T) := \{(f, \varphi) : f : T \to X \text{ a morphism},\;
  \varphi : f^* \mathcal{A} \to \mathcal{O}_T \text{ an } \mathcal{O}_T\text{-algebra map}\}.$
  Equivalently, for any \(X\)-scheme \(g : T \to X\),
  there is a natural bijection
  \[
    \Hom_X(T, \underline{\Spec}_X(\mathcal{A})) \;\;\cong\;\;
    \Hom_{\mathcal{O}_X\text{-alg}}(\mathcal{A}, g_* \mathcal{O}_T)
  \]
  (using the adjunction between \(g_*\) and \(g^*\) to convert between
  \(\mathcal{A} \to g_*\mathcal{O}_T\) and \(g^*\mathcal{A} \to \mathcal{O}_T\)).
  The bijection is natural in \(T\).
\end{theorem}

\begin{proof}
  Following Stacks tag 01LQ (proof of lemma-spec): the functor \(F\) is a Zariski sheaf
  because morphisms of schemes and \(\mathcal{O}\)-algebra maps both glue under open covers.
  Choose an affine open cover \(X = \bigcup U_i\). The subfunctor
  \(F_i \subseteq F\) of pairs \((f, \varphi)\) with \(f(T) \subseteq U_i\) is representable by
  \(\Spec(\mathcal{A}(U_i))\). Indeed, factoring \(f\) through \(U_i\) identifies
  \(f^*\mathcal{A}\) with the pullback of \(\mathcal{A}|_{U_i}\), and the affine-base
  case (\ref{mk-thm:relative_spec_affine_base}) represents the resulting algebra maps.
  Each \(F_i \hookrightarrow F\) is an open
  subfunctor (pull back \(U_i\) along \(f\)). The \(F_i\) jointly cover \(F\) since the \(U_i\)
  cover \(X\). The gluing of representable open subfunctors produces a representing scheme,
  which is canonically isomorphic to the glued scheme \(\underline{\Spec}_X(\mathcal{A})\)
  of \ref{mk-thm:relative_spec_exists}.
\end{proof}

\section{Affine base}
\label{mk-sec:relspec_affine}

\begin{theorem}[Relative spectrum over an affine base]
  \label{mk-thm:relative_spec_affine_base}
  \dcref{01LO,custom}
  \uses{mk-def:qc_sheaf_of_algebras}
  \textit{Source: \cite{stacks-project}, tag 01LO (lemma-spec-affine).}
  Let \(X = \Spec R\) be an affine scheme and let \(\mathcal{A}\) be a quasi-coherent
  \(\mathcal{O}_X\)-algebra. Set \(A := \Gamma(X, \mathcal{A})\), regarded as an
  \(R\)-algebra. Then there is a canonical isomorphism of \(X\)-schemes
  \[
    \underline{\Spec}_X(\mathcal{A}) \;\;\cong\;\; \Spec(A).
  \]
  In particular, on an affine base the relative spectrum reduces to the absolute
  spectrum of the global sections.
\end{theorem}

\begin{proof}
  Write \(X = \Spec R\) and \(A = \Gamma(X, \mathcal{A})\). Quasi-coherence of
  \(\mathcal{A}\) gives \(\mathcal{A} = \widetilde{A}\) as \(\mathcal{O}_X\)-modules
  (with the \(R\)-algebra structure on \(A\) inducing the algebra structure on
  \(\widetilde A\)). The ring map \(R \to A\) gives a morphism
  \(f_{\mathrm{univ}} : \Spec A \to \Spec R = X\) and a canonical
  \(\mathcal{O}_{\Spec A}\)-algebra map
  $\varphi_{\mathrm{univ}} : f_{\mathrm{univ}}^* \mathcal{A} = \widetilde{A \otimes_R A}
  \to \mathcal{O}_{\Spec A} = \widetilde A$ given by multiplication
  \(a \otimes a' \mapsto aa'\). We verify \((\Spec A, f_{\mathrm{univ}}, \varphi_{\mathrm{univ}})\)
  represents \(F\) on the affine base: given a test scheme \(T\) and a pair
  \((f, \varphi)\), the composite
  $A = \Gamma(X, \mathcal{A}) \xrightarrow{f^\sharp}
  \Gamma(T, f^*\mathcal{A}) \xrightarrow{\varphi} \Gamma(T, \mathcal{O}_T)$
  is a ring map \(A \to \Gamma(T, \mathcal{O}_T)\) and hence corresponds to a morphism
  \(T \to \Spec A\) by the affine universal property
  (Mathlib \texttt{Scheme.toSpec\_naturality} /
  \texttt{Scheme.Hom.toSpec}). One checks this assignment inverts the universal map
  \((f, \varphi) \mapsto a\) defined by \(f = f_{\mathrm{univ}} \circ a\),
  \(\varphi = a^* \varphi_{\mathrm{univ}}\).
\end{proof}

\section{Base change}
\label{mk-sec:relspec_basechange}

\begin{theorem}[Base change of the relative spectrum]
  \label{mk-thm:relative_spec_base_change}
  \dcref{01LS,custom}
  \uses{mk-thm:relative_spec_univ}
  \textit{Source: \cite{stacks-project}, tag 01LS (lemma-spec-base-change +
  lemma-spec-properties).}
  Let \(g : T \to X\) be a morphism of schemes and let \(\mathcal{A}\) be a quasi-coherent
  \(\mathcal{O}_X\)-algebra. Then \(g^* \mathcal{A}\) is a quasi-coherent
  \(\mathcal{O}_T\)-algebra and there is a canonical isomorphism of \(T\)-schemes
  \[
    T \times_X \underline{\Spec}_X(\mathcal{A}) \;\;\cong\;\;
    \underline{\Spec}_T(g^* \mathcal{A}).
  \]
\end{theorem}

\begin{proof}
  By \ref{mk-thm:relative_spec_univ}, \(\underline{\Spec}_X(\mathcal{A})\) represents
  the functor \(F\) sending an \(X\)-scheme \(h : S \to X\) to the set of
  \(\mathcal{O}_S\)-algebra maps \(h^*\mathcal{A} \to \mathcal{O}_S\). A morphism
  \(S \to T \times_X \underline{\Spec}_X(\mathcal{A})\) over \(T\) is the same as a pair of an
  \(X\)-morphism \(S \to \underline{\Spec}_X(\mathcal{A})\) together with a \(T\)-factorisation
  of \(S \to X\), i.e.\ a morphism \(S \to T\) whose composition with \(g\) equals
  \(S \to X\). Tracing through, this is a \(T\)-morphism \(f' : S \to T\) together with an
  \(\mathcal{O}_S\)-algebra map \((f')^*(g^*\mathcal{A}) = f^*\mathcal{A} \to \mathcal{O}_S\).
  This is precisely \(F'(S)\), the functor for \((T, g^*\mathcal{A})\), so the canonical map
  \(T \times_X \underline{\Spec}_X(\mathcal{A}) \to \underline{\Spec}_T(g^*\mathcal{A})\)
  is an isomorphism by Yoneda.
\end{proof}

\section{Pullback compatibility helpers}
\label{mk-sec:relspec_pullback_helpers}

This section records the substrate lemmas and constructions underlying the
base-change isomorphism (\ref{mk-thm:relative_spec_base_change}). Throughout,
\(g : T \to X\) is a morphism of schemes, \(\mathcal{A}\) is a quasi-coherent
\(\mathcal{O}_X\)-algebra, and \(q\) denotes the first projection of the pullback
\(T \times_X \underline{\Spec}_X(\mathcal{A})\) along the structure morphism
\(\pi : \underline{\Spec}_X(\mathcal{A}) \to X\). Write
\(\mathcal{A}_g := q_* \mathcal{O}_{T \times_X \underline{\Spec}_X(\mathcal{A})}\)
for the quasi-coherent algebra reconstructed from this affine pullback.

\begin{lemma}[Affineness of the structural pullback projection]
  \label{mk-lem:relspec_pullback_fst_isaffinehom}
  \dcref{custom}

  \uses{mk-thm:relative_spec_exists, mk-def:relspec_structure_morphism}
  The first projection
  \(q : T \times_X \underline{\Spec}_X(\mathcal{A}) \to T\) of the pullback of the
  structure morphism \(\pi\) along \(g\) is an affine morphism.
\end{lemma}

\begin{proof}
  The structure morphism \(\pi\) is affine by \ref{mk-thm:relative_spec_exists}, and
  affineness of morphisms is stable under base change; hence its pullback \(q\) is
  affine.
\end{proof}

\begin{lemma}[Quasi-coherence overlay for the pushforward along the affine projection]
  \label{mk-lem:relspec_pullback_coequifibered}
  \dcref{custom}

  \uses{mk-lem:relspec_pullback_fst_isaffinehom}
  The pushforward \(q_* \mathcal{O}_{T \times_X \underline{\Spec}_X(\mathcal{A})}\)
  along the affine projection \(q\) carries the Stacks~01LL quasi-coherence overlay:
  for every affine open \(U \subseteq T\) and section \(r \in \Gamma(T, U)\), the
  restriction of the pushforward to the basic open \(D(r) \subseteq U\) is the
  localization away from \(r\) of its sections over \(U\). This is the predicate making
  \(\mathcal{A}_g\) a well-formed quasi-coherent \(\mathcal{O}_T\)-algebra.
\end{lemma}

\begin{proof}
  Since \(q\) is affine (\ref{mk-lem:relspec_pullback_fst_isaffinehom}), the preimage
  \(q^{-1}(U)\) of an affine open \(U\) is affine, and \(q^{-1}(D(r))\) is the basic open
  of \(q^{-1}(U)\) cut out by the pulled-back section; the localization-away
  identification then follows from the affine-open localization criterion.
\end{proof}

\begin{definition}[Algebra reconstructed from the base-changed spectrum]
  \label{mk-def:relspec_qcoh_pullback}
  \dcref{custom}

  \uses{mk-def:qc_sheaf_of_algebras, mk-lem:relspec_pullback_fst_isaffinehom,
    mk-lem:relspec_pullback_coequifibered}
  Let \(g : T \to X\) be a morphism of schemes and \(\mathcal{A}\) a quasi-coherent
  \(\mathcal{O}_X\)-algebra (\ref{mk-def:qc_sheaf_of_algebras}). Define the reconstructed
  \(\mathcal{O}_T\)-algebra \(\mathcal{A}_g\) as the pushforward
  \(q_* \mathcal{O}_{T \times_X \underline{\Spec}_X(\mathcal{A})}\) along the first
  projection \(q : T \times_X \underline{\Spec}_X(\mathcal{A}) \to T\) of the pullback of
  the structure morphism \(\pi\) along \(g\). Concretely, the underlying
  \(\mathcal{O}_T\)-module is the topological pushforward \(q_* \mathcal{O}\) and the unit
  is the comorphism \(q^\sharp : \mathcal{O}_T \to q_* \mathcal{O}\); the Stacks~01LL
  quasi-coherence (coequifibered) overlay is supplied by
  \ref{mk-lem:relspec_pullback_coequifibered}, which itself rests on the affineness of \(q\)
  (\ref{mk-lem:relspec_pullback_fst_isaffinehom}). Thus \(\mathcal{A}_g\) is an object of
  \(\mathrm{QcohAlg}(T)\).
\end{definition}

\begin{definition}[Per-affine-open piece of the base-change iso]
  \label{mk-def:relspec_pullback_iso_affine_piece}
  \dcref{custom}

  \uses{mk-lem:relspec_pullback_fst_isaffinehom}
  For an affine open \(U \subseteq T\), the preimage \(q^{-1}(U)\) is affine and is
  canonically isomorphic to \(\Spec\) of the sections of \(\mathcal{A}_g\) over \(U\).
  This per-piece isomorphism feeds the global base-change
  iso construction.
\end{definition}

\begin{proof}
  As \(q\) is affine (\ref{mk-lem:relspec_pullback_fst_isaffinehom}), \(q^{-1}(U)\) is
  affine; its canonical isomorphism with \(\Spec \Gamma(q^{-1}(U), \mathcal{O})\) matches
  the sections of the pushforward \(\mathcal{A}_g\) over \(U\) by definitional
  unfolding.
\end{proof}

\begin{definition}[Pullback cocone on the relative-gluing functor]
  \label{mk-def:relspec_pullback_cocone}
  \dcref{custom}

  \uses{mk-lem:relspec_pullback_fst_isaffinehom}
  A cocone on the relative-gluing-data functor of \(\mathcal{A}_g\) whose apex is the
  pullback \(T \times_X \underline{\Spec}_X(\mathcal{A})\); its component at an affine
  open \(U\) is the canonical map \(\Spec(\mathcal{A}_g(U)) \to
  T \times_X \underline{\Spec}_X(\mathcal{A})\) into the affine piece \(q^{-1}(U)\).
\end{definition}

\begin{proof}
  The components are the affine-open inclusion maps of the pulled-back pieces; their
  compatibility (naturality) follows from the functoriality of these inclusions under
  restriction along basic opens.
\end{proof}

\begin{lemma}[Descent of the pullback cocone covers the structure morphism]
  \label{mk-lem:relspec_pullback_cocone_desc_comp_fst}
  \dcref{custom}

  \uses{mk-def:relspec_pullback_cocone, mk-lem:relspec_pullback_fst_isaffinehom}
  The colimit descent of the pullback cocone
  (\ref{mk-def:relspec_pullback_cocone}), as a morphism
  \(\underline{\Spec}_T(\mathcal{A}_g) \to
    T \times_X \underline{\Spec}_X(\mathcal{A})\),
  composed with the projection \(q\), equals the structure morphism
  \(\underline{\Spec}_T(\mathcal{A}_g) \to T\).
\end{lemma}

\begin{proof}
  By the universal property of the colimit it suffices to check the identity on each
  cocone component, where it reduces to the compatibility of the affine-piece
  inclusion with the projection \(q\).
\end{proof}

\begin{lemma}[Base-change descent is invertible]
  \label{mk-lem:relspec_pullback_iso_desc_isiso}
  \dcref{custom}

  \uses{mk-def:relspec_pullback_cocone, mk-lem:relspec_pullback_cocone_desc_comp_fst,
    mk-def:relspec_pullback_iso_affine_piece, mk-lem:relspec_pullback_fst_isaffinehom}
  The colimit descent map
  \(\underline{\Spec}_T(\mathcal{A}_g) \to
    T \times_X \underline{\Spec}_X(\mathcal{A})\)
  determined by the pullback cocone is an isomorphism of schemes.
\end{lemma}

\begin{proof}
  Being an isomorphism is local on the target, so it suffices to check it over each
  member \(q^{-1}(U)\) of the affine open cover. There the descent restricts, via
  \ref{mk-lem:relspec_pullback_cocone_desc_comp_fst} and the range identification of the
  gluing inclusions, to the per-affine-open isomorphism
  (\ref{mk-def:relspec_pullback_iso_affine_piece}).
\end{proof}

\begin{definition}[Base-change iso construction]
  \label{mk-def:relspec_pullback_iso_construction}
  \dcref{custom}

  \uses{mk-def:relspec_pullback_cocone, mk-lem:relspec_pullback_iso_desc_isiso,
    mk-lem:relspec_pullback_fst_isaffinehom}
  The canonical isomorphism of \(T\)-schemes
  \(T \times_X \underline{\Spec}_X(\mathcal{A}) \cong
    \underline{\Spec}_T(\mathcal{A}_g)\),
  defined as the inverse of the colimit descent map of the pullback cocone.
\end{definition}

\begin{proof}
  The descent map is an isomorphism by \ref{mk-lem:relspec_pullback_iso_desc_isiso}, and
  the construction takes its inverse.
\end{proof}

\begin{lemma}[Existence of the base-change isomorphism]
  \label{mk-lem:relspec_pullback_iso}
  \dcref{custom}

  \uses{mk-def:relspec_pullback_iso_construction}
  There exists a canonical isomorphism of \(T\)-schemes
  \(T \times_X \underline{\Spec}_X(\mathcal{A}) \cong
    \underline{\Spec}_T(\mathcal{A}_g)\).
\end{lemma}

\begin{proof}
  Witnessed by the explicit construction
  \ref{mk-def:relspec_pullback_iso_construction}.
\end{proof}

\section{Functoriality}
\label{mk-sec:relspec_functoriality}

\begin{theorem}[Relative Spec functor]
  \label{mk-thm:relative_spec_functorial}
  \dcref{01LR,custom}
  \uses{mk-thm:relative_spec_univ, mk-thm:relative_spec_affine_base}
  \textit{Source: \cite{stacks-project}, tag 01LR; \cite{hartshorne-ag}, II~Ex.~5.17(b).}
  The construction \(\mathcal{A} \mapsto \underline{\Spec}_X(\mathcal{A})\) assembles into
  a contravariant functor
  \[
    \underline{\Spec}_X : \mathrm{QcohAlg}(X)^{op} \longrightarrow \mathrm{AffSch}/X
  \]
  from quasi-coherent \(\mathcal{O}_X\)-algebras to affine \(X\)-schemes: a morphism of
  \(\mathcal{O}_X\)-algebras \(\psi : \mathcal{A} \to \mathcal{B}\) induces, by the universal
  property (\ref{mk-thm:relative_spec_univ}) applied to the pair
  $(\pi_{\mathcal{B}}, \psi \circ \pi_{\mathcal{B}}^*) :
  \pi_{\mathcal{B}}^*\mathcal{A} \to \pi_{\mathcal{B}}^*\mathcal{B} \to
  \mathcal{O}_{\Spec\mathcal{B}}$,
  an \(X\)-morphism
  \(\underline{\Spec}_X(\psi) : \underline{\Spec}_X(\mathcal{B}) \to \underline{\Spec}_X(\mathcal{A})\).
  Functoriality (identity, composition) is direct from the universal property. The
  pushforward \(\pi_* \mathcal{O}_{\underline{\Spec}_X(\mathcal{A})} \cong \mathcal{A}\)
  exhibits \(\underline{\Spec}_X\) as an equivalence between \(\mathrm{QcohAlg}(X)^{op}\) and
  the full subcategory of affine \(X\)-schemes, with inverse \(f \mapsto f_* \mathcal{O}_Y\)
  (Stacks lemma-spec-properties part~(3)).
\end{theorem}

\begin{proof}
  Functoriality follows directly from \ref{mk-thm:relative_spec_univ}: for
  \(\psi : \mathcal{A} \to \mathcal{B}\), the pair
  \((\pi_{\mathcal{B}} : \underline{\Spec}_X(\mathcal{B}) \to X, \;
  \pi_{\mathcal{B}}^*(\psi))\) defines an \(X\)-morphism
  \(\underline{\Spec}_X(\mathcal{B}) \to \underline{\Spec}_X(\mathcal{A})\) via the
  representing universal pair. Identity and composition follow from naturality of the
  representing bijection in \(\mathcal{A}\). For the equivalence claim, \(\underline{\Spec}_X\)
  is fully faithful because the isomorphism
  \(\pi_* \mathcal{O}_{\underline{\Spec}_X(\mathcal{A})} \cong \mathcal{A}\)
  (locally identifying with \(\Gamma(\Spec(\mathcal{A}(U)),\mathcal{O}) = \mathcal{A}(U)\)
  via \ref{mk-thm:relative_spec_affine_base}) reconstructs \(\mathcal{A}\) from
  \(\underline{\Spec}_X(\mathcal{A})\). Essential surjectivity onto affine \(X\)-schemes uses
  the canonical morphism \(X \to \underline{\Spec}_X(\pi_*\mathcal{O}_X)\) of Stacks
  lemma-canonical-morphism; on an affine \(X\)-scheme \(\pi : Y \to X\), this canonical
  morphism is an isomorphism (affine-locally it is the identity
  \(\Spec\Gamma(\pi^{-1}(U), \mathcal{O}_Y) \to \Spec\Gamma(\pi^{-1}(U), \mathcal{O}_Y)\)).
\end{proof}

\section{Construction by affine gluing}
The relative spectrum is determined by the affine schemes
\(\Spec(\mathcal{A}(U))\) over the affine opens \(U \subseteq X\). Quasi-coherence
identifies their restrictions on overlaps, and the resulting cocycle glues both the
schemes and their structure morphisms. The universal property shows that this glued
object is independent of the chosen affine cover.

\section{Further properties}
\label{mk-sec:relspec_out_of_scope}

Additional hypotheses on \(\mathcal{A}\), such as finite presentation, flatness, or
smoothness, impose corresponding finiteness and geometric properties on
\(\underline{\Spec}_X(\mathcal{A}) \to X\). The relative Picard construction uses the
functoriality and base-change results above together with such properties in its own
geometric setting.
